\begin{Example}
\label{ex:electrostatic-surface-charge-weak-solution}
一系列重要的经典 Dirichlet 问题例子来自静电学. 不涉及任何物理细节, 假设希望知道区域 $\Omega\subset\mathbb{R}^3$ 中一个称为\emph{静电势}, 或更常称为\emph{电压}的标量 $u$. 假设区域的边界接地; 即
\begin{equation}\label{eq:electrostatic-grounded-dirichlet-boundary}
u(\bm{x})=0\quad\text{for }\bm{x}\in\partial\Omega.
\end{equation}
在 $\Omega$ 内有电荷分布 $\rho:\Omega\to\mathbb{R}$, 这些电荷通过公式
\begin{equation}\label{eq:electrostatic-charge-poisson-equation}
-\Delta u=\rho.
\end{equation}
``产生\textquotesingle\textquotesingle\ 电压. 这类问题的求解是经典位势理论的研究内容, 经典理论的许多方法 (特征函数展开, Green 函数) 也包含在现代理论中. 然而, 只有当寻找经典解 ($\rho\in C_b(\Omega)$) 或强解 ($\rho\in L^2(\Omega)$ 或分段连续) 时, 经典理论与现代理论才采用相同的方法.

对于分布在曲面上的电荷, 经典理论与现代理论采取截然不同的方法; 即 $S\subset\Omega$ 为光滑曲面, 并在其上定义电荷分布 $\omega:S\to\mathbb{R}$ 的情形. (假设 $\omega\in C_b(S)$ 或 $\omega\in L^2(S)$.) 当然, 由于 $\omega$ 只定义在一个曲面上 (对学过测度论的读者而言, 这是零测集), 微分方程 $-\Delta u=\omega$ 无论按经典意义还是按 $L^2(\Omega)$ 中序列等价类的等同来理解, 都没有意义. 在现代理论中, 情况十分清楚: 虽然 $\omega$ 不属于 $L^2(\Omega)$, 却可以用公式
\begin{equation}\label{eq:surface-charge-negative-sobolev-functional}
(\omega,\phi)=\int_S\omega(\bm{x})\phi(\bm{x})\,\mathrm{d}a(\bm{x}),
\end{equation}
对每个 $\phi\in H_0^1(\Omega)$ 定义一个很好的 $H^{-1}(\Omega)$ 中的泛函. (回想迹定理表明 $\phi\in L^2(S)$.) 这里 $\mathrm{d}a(\bm{x})$ 表示 $\bm{x}\in S$ 处的面积微元. 因而现代理论只要求寻找变分问题
\begin{equation}\label{eq:surface-charge-electrostatic-variational-problem}
B[v,u]:=\int_\Omega\nabla v\cdot\nabla u\,\mathrm{d}\bm{x}=(\omega,v)
\end{equation}
对所有 $v\in H_0^1(\Omega)$ 成立的一个 (弱) 解 $u\in H_0^1(\Omega)$. 如下文所示, 该问题适定. 经典理论利用单层势与双层势理论来求解这一问题 (及其他一些类似问题): 其本质是由奇异曲面积分定义的积分算子. 正如经常出现的情形一样, 经典理论缺乏现代理论在概念上的许多统一性, 但在特殊情形 (尽管往往是最重要的情形) 中提供了详细得多的信息. 本书不讨论经典位势理论的结果, 但建议读者阅读 O.D. Kellogg 的 \emph{Foundations of Potential Theory} \cite{Ke}, 作为进一步了解这一主题的良好起点.
\end{Example}
